\documentclass[10pt]{article}
\usepackage{graphicx} 

\usepackage{C:/Users/milen/OneDrive/Desktop/usefullstuff}

\usepackage{titlesec}
% \hypersetup{hidelinks}
\hypersetup{
    colorlinks = true,
    urlcolor  = {blue!50!black}, 
}
% \usepackage[left=1in, right=1in, top=1in, bottom=1in]{geometry}
\usepackage{fontawesome5}
\geometry{margin=1in}

\titleformat{\section}{\large\bfseries}{}{0em}{}[\titlerule]

\begin{document}

\pagestyle{empty}
\begin{center}
    {\textbf{\huge Milène Mandeville}}\\\Large{Applied Mathematics Student}\vspace{3mm}
\end{center}
\begin{center}
    \faIcon{phone} \href{https://wa.me/33668298113}{+33 6 68 29 81 13} \hspace{1cm}
    \faIcon{envelope} \href{mailto:milenemandeville@yahoo.com}{milenemandeville@yahoo.com} \hspace{1cm}
    \faIcon{linkedin} \href{https://www.linkedin.com/in/milène-mandeville}{milènemandeville}\vspace{3mm}
\end{center}
\begin{minipage}[t]{0.3\linewidth}
\section*{Languages}
\begin{itemize}
    \item French (native)
    \item English (C1)
    \item Dutch (B1)
    \item German (B1)
\end{itemize}
\end{minipage}\hfill
\begin{minipage}[t]{0.3\linewidth}
\section*{Programming skills}
\begin{itemize}
    \item Python (advanced)
    \item Matlab (intermediate)
    \item Julia (beginner)
\end{itemize}
\end{minipage}\hfill
\begin{minipage}[t]{0.3\linewidth}
\section*{Diplomas}
\begin{itemize}
    \item BSc in Applied \\ Mathematics (2024)
    \item French Baccalaureate in Science -- honours (2020)
\end{itemize}
\end{minipage}\hfill 

\section*{Research}
\begin{itemize}
    \item \textbf{Master Thesis}: University of Groningen (2026, ongoing) \begin{itemize}
        \item \textbf{(Working) Title}: Extension of a 1D model for large eddy simulation in Burger’s turbulence
        \item \textbf{Supervisors} prof. dr. ir. R.W.C.P Verstappen and dr. ir. R. Luppes 
        \item \textbf{Description}: The aim of this work is to extend an
        existing model (Verstappen, 2026) for large eddy simulation in Burger’s turbulence. This
        one dimensional
        model relies on a uniform mesh with periodic boundary conditions, and we
        wish to generalise it to a three dimensional triangular grid. 
    \end{itemize}
    \item \textbf{Internship}: Von Karman Insitute for Fluid Dynamics (2026) \begin{itemize}
        \item \textbf{Title}: \href{https://fse.studenttheses.ub.rug.nl/37239/1/mAppM2026MandevilleMA.pdf}{Towards Efficient Galerkin Projection in
        Lagrangian ROMs for Sloshing}
        \item \textbf{Supervisors}: prof. M. A. Mendez, ir. S. Ahizi,  dr. J. Koellermeier, and dr. ir. R.
        Luppes,
        \item \textbf{Description}: In this paper I derive a reduced order model for
        smoothed particle hydrodynamics, by rewriting the standard neighbour
        sums as dense matrix operations, and using proper orthorgonal
        decomposition and Galerkin projection to obtain a linear time invariant
        model. However, the kernel transformations inherent to particle dynamics
        introduce a nonlinear state dependent term. To this end, I evaluate
        three methods: linearisation, random Fourier features, and a singular
        value decomposition based interpolation. 
    \end{itemize}
    \item \textbf{Bachelor Thesis}: University of Groningen (2023) \begin{itemize}
        \item \textbf{Title}: \href{https://fse.studenttheses.ub.rug.nl/30589/1/bAppM_2023_MandevilleMA.pdf}{Average Surface Roughness for 2-Dimensional Atmospheric Flow}
        \item \textbf{Supervisors}: dr. ir. R. Luppes and dr. A. E. Sterk 
        \item \textbf{Description}: In this paper I present a study of four different methods of computing an
        average (effective) roughness length (following André and Blondin (1986), Taylor
        (1989), de Vries et. al. (2003), and taking the average of the
        logarithmic law).
    \end{itemize}
\end{itemize}

\section*{Professional experience}
\begin{itemize}
\item \textbf{Acadomia (tutoring company): online teacher} (2024--present)
\begin{itemize}
    \item Teaching mathematics to middle school and high school students in French
\end{itemize}
\item \textbf{University of Groningen: teaching assistant} (2023--present)\\ 
    For the following Bachelor courses, I gave tutorials and/or computer labs. I
    also graded homeworks and exams, and evaluated students
    with oral examinations for the courses with computer labs. \begin{itemize}
    \item Mathematical Tools for BME (2025) -- C. Karakurt, PhD
    \item Project Mathematical Modelling (2025) -- prof. dr. ir. R. Verstappen
    \item Computational Methods of Science (2024--2026) -- prof. dr. ir. R. Verstappen, dr. ir. R. Luppes
    \item Mathematical Modelling (2023--2026) -- dr. S. Trenn, prof dr. ir. R. Verstappen
    \item Advanced Systems Theory (2023) -- prof. M.K. Camlibel
\end{itemize}
    % \item \textbf{Puteaux townhall Public Library (Paris La Défense suburb): summer job} (2023) \begin{itemize}
%     \item Welcomed and helped visitors
%     \item Sorted books using Orphée software
% \end{itemize}
% \item \textbf{Penoyre \& Prasad Architectes: intern} (2017)
% \begin{itemize}
%     \item Observational internship
% \end{itemize}
\end{itemize}


\section*{Other experience}
\begin{itemize}
    \item \textbf{Mathematica committee of the FMF (study association): secretary} (2023--present)\begin{itemize}
    \item Writing minutes summarising the meetings 
    \item Writing emails inviting professors to give a talk 
    \item Organising the annual integration competition and a series of informal lectures
\end{itemize}
\end{itemize}

\section*{Education}
\begin{itemize}
    \item \textbf{MSc Applied Mathematics}, Rijksuniversiteit Groningen (2024--present)
    \item \textbf{BSc Applied Mathematics}, Rijksuniversiteit Groningen (2020-2024)
    % \item Agora High School, Puteaux France (2017--2020)
\end{itemize}

\qed

\end{document}
